\section{Chaîne audio filaire}
\label{sec:audio_filaire}

Les deux sorties du lecteur, la prise casque et le Bluetooth, sont exposées au reste du firmware derrière une interface unique, le \gls{sink}, qui se réduit à ouvrir une sortie dans un format donné, y écrire un bloc d'échantillons, puis la fermer. Le transport de lecture ignore ainsi vers quelle sortie part le son, et la bascule de l'une à l'autre en cours de lecture n'impose aucune duplication du pipeline. Cette section décrit l'implémentation filaire de ce sink, la variante Bluetooth étant traitée à la section~\ref{sec:soft_bluetooth}.

\subsection{Bus et format des échantillons}

L'horloge maître (\gls{mclk}) du bus \gls{i2s} n'est pas câblée sur cette carte, ce qui économise une broche et une piste. Le bus se réduit à trois signaux et les conséquences de ce choix sur l'arbre d'horloge sont traitées à la section~\ref{sec:horloges_dac}~\cite{TI_PCM5242}.

Le pipeline manipule des échantillons de 32~bits par canal, justifiés à gauche et entrelacés, quelle que soit la résolution du fichier source. Ce choix est imposé par ESP-IDF~6, qui n'élargit plus automatiquement les échantillons à la largeur du \gls{slot} matériel~\cite{ESP32_IDF_migration_v6}. Un flux de 16~bits transmis tel quel remplirait un même slot avec deux échantillons consécutifs et serait lu au double de la vitesse attendue. L'élargissement est donc réalisé par les décodeurs, en amont du transfert I\textsuperscript{2}S.

La largeur des slots reste en revanche libre, le convertisseur acceptant une horloge de bit comprise entre 32 et 256 fois la fréquence d'échantillonnage~\cite{TI_PCM5242}. Elle est donc fixée à 32~bits de largeur en toutes circonstances. Les mots sont ainsi recopiés sans reformatage et l'horloge de bit vaut toujours 64 fois la fréquence d'échantillonnage, soit deux slots de 32~bits par trame, ce qui autorise une table de configuration unique du convertisseur. Le prix à payer se limite à de la bande passante, l'horloge de bit culminant à 12,288~MHz.

\subsection{Génération des horloges}
\label{sec:horloges_dac}

L'ESP32 génère par défaut l'horloge de bit via un diviseur fractionnaire de 160~MHz, introduisant une irrégularité des fronts déconseillée pour l'audio~\cite{ESP32_IDF_i2s}. La documentation recommande l'usage de la boucle analogique dédiée (\gls{apll})~\cite{ESP32_Technical_Reference_Manual}. Le choix est crucial ici, le convertisseur asservissant sa propre boucle sur BCK en l'absence de \gls{mclk} externe. Toute irrégularité des fronts se propagerait donc jusqu'à son modulateur interne. L'\gls{apll} est retenue comme horloge source du périphérique I\textsuperscript{2}S, dont sont dérivées BCK et WS, son surcoût énergétique n'intervenant que durant la lecture.

\fig[H, width=0.85\linewidth]{Interface I\textsuperscript{2}S à trois fils entre le microcontrôleur et le DAC. Diagramme créé avec l'aide de Claude Opus 4.8.}{Lien_I2S.drawio}

Le convertisseur dispose d'un mode de détection automatique des horloges mais qui ne fonctionne pas lorsque le circuit est piloté par I\textsuperscript{2}C~\cite{TI_PCM5242}. L'arbre d'horloge est donc programmé manuellement. Le firmware place le convertisseur en veille, lui ordonne d'ignorer cette broche, bascule la référence de sa boucle sur l'horloge de bit, écrit les coefficients puis le réveille, chaque écriture étant relue et vérifiée.

\fig[H, width=0.7\textwidth]{Arbre d'horloge interne du PCM5242. Diagramme créé avec l'aide de Claude Opus 4.8.}{Arbre_horloge_PCM5242.drawio}

L'oscillateur ne prend que deux valeurs, 90,3168 et 98,304~MHz, soit 2048 fois 44,1 et 48~kHz. Chacune est un multiple entier de toutes les fréquences de sa famille, condition pour que les divisions successives de l'arbre retombent exactement sur la fréquence demandée. Aucune autre valeur ne satisfait ce critère dans la fenêtre de 64 à 100~MHz imposée par la fiche technique~\cite{TI_PCM5242}.

Les coefficients de la boucle s'en déduisent. L'oscillateur vaut l'horloge de bit multipliée par $J \times R$, d'où $J \times R = f_{\text{VCO}} / (64 f_S)$. L'oscillateur restant fixe au sein d'une famille, ce rapport est divisé par deux chaque fois que la fréquence d'échantillonnage double, de 32 à 44,1~kHz jusqu'à 8 à 192~kHz. La table complète figure en annexe~\ref{ann:horloges}.

La plage supportée s'étend de 16 à 192~kHz. La limite supérieure répond au cahier des charges tandis que la limite inférieure découle d'une contrainte technique, le convertisseur exige une référence d'au moins 1~MHz, ce qui interdit toute fréquence inférieure à 15,6~kHz avec le rapport retenu~\cite{TI_PCM5242}. Toute demande hors plage est explicitement refusée pour éviter un dysfonctionnement silencieux.

\subsection{Tampons de transfert}

Le contrôleur \gls{dma} puise les échantillons dans un anneau de tampons que la tâche audio remplit à l'avance. Cet anneau absorbe les retards de cette tâche, qui dépend de la carte SD dont le temps de réponse n'est pas borné. La configuration par défaut de six descripteurs de 240~trames ne couvre que 7,5~ms à 192~kHz et a provoqué des coupures lors des premiers tests sur matériel~\cite{ESP32_IDF_i2s}. L'anneau a donc été porté à douze descripteurs de 511~trames. Cette valeur de 511 est un plafond matériel, car un descripteur ne peut décrire plus de 4092~octets et une trame stéréo de 32~bits par canal en pèse 8~\cite{ESP32_Technical_Reference_Manual}. La réserve atteint alors $t_{192} = \frac{12 \times 511}{192\,000} \approx 32~\text{ms}$ pour 49\,056~octets, soit environ 48~kio de RAM interne. Le \gls{dma} n'adressant pas la PSRAM~\cite{ESP32_IDF_guide}, cette dépense n'a été possible que grâce au déport du tampon d'affichage en mémoire externe, selon l'arbitrage de la section~\ref{sec:archi}. Un anneau plus profond retarderait en revanche l'effet d'un arrêt ou d'un changement de piste, car les trames déjà écrites sont jouées jusqu'au bout.

\subsection{Volume et séquencement}

La position du potentiomètre est lue par l'ADC de l'ESP32. il ne sert que de capteur et n'est pas inséré dans le chemin du signal. Sa course est projetée linéairement sur une plage en décibels plutôt qu'en amplitude, car les registres de volume du DAC s'expriment eux-mêmes en dB. Une variation linéaire de ces registres produit donc une perception d'intensité sonore linéaire à l'oreille.

Le gain logiciel maximal a été fixé à 0~dB lors de l'implémentation, bien que le DAC autorise un gain positif jusqu'à +24~dB \cite{TI_PCM5242}. Cette marge n'est volontairement pas exploitée car l'amplificateur, avec un gain fixe de 0,5, utilise déjà toute sa plage dynamique. Le DAC délivre 4,2~Vrms en différentiel, ramenés à 2,1~Vrms après amplification, une valeur proche de son excursion maximale de 2,15~Vrms sous 3,3~V \cite{analog_MAX97220A}. Un gain logiciel supplémentaire risquerait donc de saturer l'amplificateur. Cette limite pourra être révisée si les tests révèlent une saturation imprévue.

L'atténuation est appliquée directement dans le cœur numérique du DAC, ce qui l'exempte de la latence de l'anneau de tampons. Le service de volume est le seul à écrire ce registre et n'émet une transaction que lorsque sa valeur change réellement.

L'ordre d'activation des étages conditionne l'absence de bruit parasite au casque. La séquence programme la fréquence d'échantillonnage, puis le volume, sort le convertisseur de veille et n'active l'amplificateur qu'en dernier, celui-ci jouant le rôle d'une porte fermée pendant l'établissement de la chaîne amont. L'extinction suit l'ordre inverse, déclenchée par un point d'accroche que le service d'énergie appelle avant de relâcher l'alimentation, ce qui lui évite toute connaissance de la chaîne audio et préserve la règle de dépendance descendante. Le même encadrement protège les changements de fréquence d'échantillonnage en cours de session, la reprogrammation de la boucle faisant nécessairement perdre son verrouillage au convertisseur.